\chapter{Correspondencia con las lecturas}
Las veintisiete lecturas originales disponibles del cuarto curso se integran en seis partes y dieciséis capítulos. Se desarrollan sus \textbf{81 ejercicios numerados}, con todos sus incisos. La lectura 24 es una reconstrucción editorial independiente y añade \textbf{ocho ejercicios}; se distingue de los originales tanto aquí como en el capítulo~\ref{cap:v4-13}. Los diez ejercicios de ampliación completan un total de \textbf{99 ejercicios con solución}.

La tabla cuenta ejercicios, no incisos ni verificaciones incluidas dentro de las pruebas. Las lecturas 5, 7, 12, 14, 18, 19 y 26--28 no contienen listas numeradas de ejercicios; sus preguntas y demostraciones propuestas dentro de la teoría se desarrollan en el texto. Las herramientas de convolución e identidades aproximadas empleadas por la reconstrucción 24 se estudian previamente en el capítulo~\ref{cap:v4-12}. El ejercicio~\ref{ej:v4-l25-2} utiliza la dualidad y la extensión de caracteres demostradas en el capítulo~\ref{cap:v4-16}; la dependencia se indica en su solución.

Los encabezados conservan la lectura y el número de origen; la numeración principal sigue cada capítulo. Las correcciones se explican en el texto. El PDF de la lectura 19 lleva por error el encabezado «Lectura 18».

\footnotesize
\begin{longtable}{@{}r>{\raggedright\arraybackslash}p{.51\textwidth}rrr@{}}
\toprule Lectura & Tema & Cap. & Ej. & Pág.\\\midrule
\endfirsthead
\toprule Lectura & Tema & Cap. & Ej. & Pág.\\\midrule
\endhead
\bottomrule\endfoot
1 & Sumación de Cesàro & \ref{cap:v4-01} & 8 & \pageref{lectura:v4-01}\\
2 & Integral de Lebesgue en la recta & \ref{cap:v4-02} & 4 & \pageref{lectura:v4-02}\\
3 & Espacios periódicos y familias trigonométricas & \ref{cap:v4-02} & 2 & \pageref{lectura:v4-03}\\
4 & Series de Fourier trigonométricas & \ref{cap:v4-03} & 4 & \pageref{lectura:v4-04}\\
5 & Derivación, integración y sumación de Fejér & \ref{cap:v4-03} & 0 & \pageref{lectura:v4-05}\\
6 & Fourier en $L^2$ y Parseval & \ref{cap:v4-04} & 7 & \pageref{lectura:v4-06}\\
7 & Localización, Dini y puntos de Lebesgue & \ref{cap:v4-04} & 0 & \pageref{lectura:v4-07}\\
8 & Convolución periódica & \ref{cap:v4-05} & 4 & \pageref{lectura:v4-08}\\
9 & Convolución en la recta, extendida a $\R^d$ & \ref{cap:v4-06} & 5 & \pageref{lectura:v4-09}\\
10 & Transformada de Fourier en la recta & \ref{cap:v4-07} & 3 & \pageref{lectura:v4-10}\\
11 & Inversión y funciones de Schwartz & \ref{cap:v4-07} & 5 & \pageref{lectura:v4-11}\\
12 & Caracteres de $L^1(\R)$ & \ref{cap:v4-07} & 0 & \pageref{lectura:v4-12}\\
13 & Plancherel y cálculos de integrales & \ref{cap:v4-08} & 4 & \pageref{lectura:v4-13}\\
14 & Subespacios invariantes por traslaciones & \ref{cap:v4-08} & 0 & \pageref{lectura:v4-14}\\
15 & Grupos topológicos & \ref{cap:v4-09} & 4 & \pageref{lectura:v4-15}\\
16 & Cocientes y estructura topológica & \ref{cap:v4-09} & 9 & \pageref{lectura:v4-16}\\
17 & Preparación para la medida de Haar & \ref{cap:v4-10} & 4 & \pageref{lectura:v4-17}\\
18 & Construcción de Haar & \ref{cap:v4-10} & 0 & \pageref{lectura:v4-18}\\
19 & Unicidad de Haar & \ref{cap:v4-10} & 0 & \pageref{lectura:v4-19}\\
20 & Función modular & \ref{cap:v4-11} & 3 & \pageref{lectura:v4-20}\\
21 & Integración en grupos cociente & \ref{cap:v4-11} & 3 & \pageref{lectura:v4-21}\\
22 & Álgebra de convolución & \ref{cap:v4-12} & 5 & \pageref{lectura:v4-22}\\
23 & Representaciones unitarias y Schur & \ref{cap:v4-13} & 4 & \pageref{lectura:v4-23}\\
24 & Representaciones integradas (reconstrucción) & \ref{cap:v4-13} & 8 & \pageref{lectura:v4-24}\\
25 & Caracteres y grupo dual & \ref{cap:v4-14} & 3 & \pageref{lectura:v4-25}\\
26 & Fourier en grupos LCA & \ref{cap:v4-15} & 0 & \pageref{lectura:v4-26}\\
27 & Álgebra reducida y espectro & \ref{cap:v4-16} & 0 & \pageref{lectura:v4-27}\\
28 & Regularización, positividad e inversión & \ref{cap:v4-16} & 0 & \pageref{lectura:v4-28}\\
\end{longtable}
\normalsize
